% ECONOMICS_PROBLEM: {"id": "H63-443", "title": "Optimal sovereign-debt maturity", "jel_code": "H63", "source_id": "OP-0257"}
% !TeX program = xelatex
\documentclass[11pt,a4paper]{article}
\usepackage[margin=23mm]{geometry}
\usepackage{fontspec}
\setmainfont{Latin Modern Roman}
\usepackage{amsmath,amssymb,mathtools}
\usepackage{xcolor,enumitem,booktabs,longtable,array,xurl,hyperref}
\definecolor{muted}{HTML}{53636C}
\hypersetup{colorlinks=true,urlcolor=blue,linkcolor=blue}
\setlength{\parindent}{0pt}
\setlength{\parskip}{5pt}
\newcommand{\R}{\mathbb R}
\newcommand{\E}{\mathbb E}
\newcommand{\N}{\mathbb N}
\newcommand{\ind}{\mathbf 1}
\newcommand{\Var}{\operatorname{Var}}
\newcommand{\Cov}{\operatorname{Cov}}
\newcommand{\supp}{\operatorname{supp}}
\DeclareMathOperator*{\argmin}{arg\,min}
\DeclareMathOperator*{\argmax}{arg\,max}
\newcommand{\entrymeta}[3]{{\sffamily\small\color{muted}\textbf{\texttt{#1}}\quad|\quad #2\par #3\par}}
\newcommand{\Problemref}[1]{\texttt{#1}}
\newcommand{\JELref}[2]{\texttt{#2} (JEL \texttt{#1})}
\begin{document}
\section*{Optimal sovereign-debt maturity}
\textbf{Problem ID:} \texttt{H63-443}\quad \textbf{JEL:} \texttt{H63}\par
% BEGIN ECONOMICS STATEMENT
\label{problem:OP-0257}
{\sffamily\small\color{muted}Original subject group: Fiscal policy and sovereign debt\par}
\label{definitions:problem:30007538}
\textbf{Audit classification:} Explicit published open question. Rollover-risk measurement and maturity evidence gaps are explicit; the dynamic optimal-control problem is editorial. Verification concerns the cited version, not first priority or proof that this exact editorial specification remains unsolved on October 8, 2026.
\subsection*{Definitions and precise research target}
\paragraph{Editorial mathematical specification.} Consider a sovereign with stochastic output \(y_t\), inherited short- and long-term debt \((b_t^S,b_t^L)\), and a specified cost of default. Short debt matures next period; long debt has declared coupon and decay rate. Investors price both claims using a stated discount factor, recovery rule, and information structure. Government decisions jointly determine primary surpluses, default, and new issuance; equilibrium prices \(q^S,q^L\) therefore depend on the policy and debt state.
Let \(m_t\in[0,1]\) be the fraction of new funding raised at long maturity. For a fixed initial state, determine a policy \(m^*(y,b^S,b^L)\) maximizing
\[
E\sum_{t\ge0}\beta^t
[u(c_t)-\mathcal D_t],
\]
where \(c_t\) is resident consumption after real default expenses, \(\beta\in(0,1)\) the resident discount factor, and \(\mathcal D_t\) additional nonconsumption disruption loss measured in utility units. Incorporate term premia, refinancing needs, and the opportunity to adjust fiscal policy after future information. Establish comparative statics and quantitatively disciplined welfare gains relative to fixed maturity shares. Separate fundamental insolvency from rollover failure through explicit equilibria or identification restrictions. Validate model-implied maturity and spread responses rather than fitting average debt composition alone. The target is optimal maturity in a declared endogenous-price environment; existing models and practical maturity strategies already solve important special cases, while empirical measurement and generalization remain unfinished.

Write net new one-period issuance as \(n_t^S\), long issuance as \(n_t^L\), and long debt survival as \(1-\delta\), \(0<\delta<1\), with coupon \(\kappa_L\ge0\). Without default, \(b_{t+1}^S=n_t^S\) and \(b_{t+1}^L=(1-\delta)b_t^L+n_t^L\). Specify the budget \(c_t=y_t-\text{real government purchases and default resource costs}-b_t^S-(\delta+\kappa_L)b_t^L+q_t^Sn_t^S+q_t^Ln_t^L\), adjusted for the declared output and recovery effects of default. Define the funding maturity share as \(m_t=q_t^Ln_t^L/(q_t^Sn_t^S+q_t^Ln_t^L)\) only when gross new funding is positive and both issuances are nonnegative; permit buybacks through separate controls. A Markov policy needs every price-relevant state, including a refinancing/sentiment state if present. A maturity rule alone cannot maximize welfare unless fiscal and default policies are optimized jointly or held explicitly fixed.
\subsection*{Variants and assumption changes}
The following are \emph{editorial variants}, not additional published open conjectures.
\begin{enumerate}
\item \emph{No self-fulfilling rollover risk.} Impose a unique equilibrium with repayment risk determined solely by fundamentals. Characterize optimal issuance maturity given term premia and endogenous default, then compare with an exogenous fixed-price benchmark.
\item \emph{Sentiment and liquidity buffer.} Add a refinancing sunspot and a reserve buffer as an additional control. Quantify how long maturity and liquid reserves substitute or complement one another in avoiding crisis under a stated equilibrium selection.
\end{enumerate}
\subsection*{References and source audit}
\begin{itemize}
\item Kris James Mitchener; Christoph Trebesch. \emph{Sovereign Debt in the 21st Century}. 2021. Locator: Section6, printed p.35, paragraphs on empirical rollover risk and maturity; corrected from Section7. Primary text checked. \url{https://www.ifo.de/DocDL/cesifo1_wp8959.pdf}.
\end{itemize}
Rollover-risk measurement and maturity evidence gaps are explicit; the dynamic optimal-control problem is editorial.
\begin{enumerate}\raggedright
\item\hypertarget{definitions:source:30007538:1}{} Kris James Mitchener; Christoph Trebesch. 2021. \emph{Sovereign Debt in the 21st Century}. Section6, printed p.35, paragraphs on empirical rollover risk and maturity; corrected from Section7.
\url{https://www.ifo.de/DocDL/cesifo1_wp8959.pdf}.
\end{enumerate}

\subsection*{Common conventions: Source mathematical conventions}

These conventions apply to each entry unless explicitly replaced.

\textbf{Models and identification.} A primitive model \(m\) specifies all probability laws, feasible choices, information, equilibrium and measurement rules used by its target. The admissible class \(\mathcal M\) is fixed independently of outcomes. If \(P_O(m)\) is its observable probability law and \(\tau(m)\) the target, its identified set at observable law \(P\) is
\[
 \mathcal I_\tau(P)=\{\tau(m):m\in\mathcal M,\ P_O(m)=P\}.
\]
Point identification means that this set is a singleton. An empty set signals incompatibility of data and assumptions. Coordinate bounds need not describe the joint set. Bounds are sharp only if every retained target is attainable or arbitrarily approachable by compatible models. Finite bounds require explicit restrictions on unobserved quantities; unrestricted confounding, hidden wealth, extrapolation or arbitrary equilibrium selection can yield uninformative sets.

\textbf{Causal interventions.} A potential outcome \(Y_i(p)\) is the outcome under a complete policy regime \(p\), including timing, financing, information and market spillovers. The target population and units are fixed before treatment. With interference, use the complete assignment vector or a justified exposure mapping; individual no-interference notation alone cannot identify an economy-wide intervention. A difference of mean potential outcomes does not require the cross-world joint distribution. Conditioning on post-treatment migration, survival, enrollment or employment changes the estimand and needs separate justification.

\textbf{Statistical statements.} An estimator, confidence set, forecast or test specifies a sampling law, model class and loss/coverage criterion. Uniform \((1-\alpha)\)-coverage means \(\inf_{m\in\mathcal M}\Pr_m\{\tau(m)\in C_n\}\geq1-\alpha\), or an explicitly asymptotic version. The whole parameter space trivially has coverage, so informativeness requires a stated width, risk or power objective. A forecasting improvement is relative to a named benchmark, information set and evaluation population.

\textbf{Equilibrium and optimization.} An equilibrium comprises feasible plans, optimality, beliefs where needed, budget constraints and all market/resource clearing. The primitives or a stated benchmark define each correspondence \(\mathcal E(p;m)\). Nonemptiness is assumed or proved, never inferred just from compact policy sets. A selected-equilibrium target uses a declared selection rule; a robust target quantifies over all permitted equilibria. Use suprema or infima unless attainment is established. An \(\varepsilon\)-optimal policy is feasible and within \(\varepsilon\) of the finite optimum value. Report an empty feasible set separately.

\textbf{Welfare and accounting.} Normative utility, interpersonal normalization, social weights, numeraire, discounting and the use of public receipts are primitives. Behavioral utility need not coincide with the welfare criterion. Household welfare includes ownership income and rents once; adding profits again to compensating variation generally double-counts them. A monetary equivalent is defined by an explicit transfer and price convention, with monotonicity and range conditions. Average effects, mechanism contributions, transfers and welfare losses are different targets. Infinite sums require convergence and dynamic feasible plans require terminal or no-Ponzi conditions.

\textbf{Source versions.} A working-paper date, revision date and journal publication year are distinguished. A DOI for a journal version is not automatically the DOI of an earlier manuscript. Locators refer to the stated version and printed pages where specified. A metadata-only check does not verify a claim about a full-text conclusion. Every entry's source subsection narrows the provenance claim accordingly.
% END ECONOMICS STATEMENT
\end{document}
